\documentclass[10pt,a4paper]{amsart}
\usepackage{amsmath,amssymb,amsthm,graphicx}
\usepackage{verbatim}
\newenvironment{reference}{\comment}{\endcomment}
%\newenvironment{reference}{}{}
\newenvironment{slogan}{\comment}{\endcomment}
\newenvironment{history}{\comment}{\endcomment}

% For commutative diagrams we use Xy-pic
\usepackage[all]{xy}

% We use 2cell for 2-commutative diagrams.
\xyoption{2cell}
% Theorem environments.
%
\theoremstyle{plain}
\newtheorem{theorem}[subsection]{Theorem}
\newtheorem{proposition}[subsection]{Proposition}
\newtheorem{lemma}[subsection]{Lemma}

\theoremstyle{definition}
\newtheorem{definition}[subsection]{Definition}
\newtheorem{example}[subsection]{Example}
\newtheorem{exercise}[subsection]{Exercise}
\newtheorem{situation}[subsection]{Situation}

\theoremstyle{remark}
\newtheorem{remark}[subsection]{Remark}
\newtheorem{remarks}[subsection]{Remarks}

\numberwithin{equation}{subsection}

% Macros
%
\def\lim{\mathop{\mathrm{lim}}\nolimits}
\def\colim{\mathop{\mathrm{colim}}\nolimits}
\def\Spec{\mathop{\mathrm{Spec}}}
\def\Hom{\mathop{\mathrm{Hom}}\nolimits}
\def\Ext{\mathop{\mathrm{Ext}}\nolimits}
\def\SheafHom{\mathop{\mathcal{H}\!\mathit{om}}\nolimits}
\def\SheafExt{\mathop{\mathcal{E}\!\mathit{xt}}\nolimits}
\def\Sch{\mathit{Sch}}
\def\Mor{\mathop{\mathrm{Mor}}\nolimits}
\def\Ob{\mathop{\mathrm{Ob}}\nolimits}
\def\Sh{\mathop{\mathit{Sh}}\nolimits}
\def\NL{\mathop{N\!L}\nolimits}
\def\CH{\mathop{\mathrm{CH}}\nolimits}
\def\proetale{{pro\text{-}\acute{e}tale}}
\def\etale{{\acute{e}tale}}
\def\QCoh{\mathit{QCoh}}
\def\Ker{\mathop{\mathrm{Ker}}}
\def\Im{\mathop{\mathrm{Im}}}
\def\Coker{\mathop{\mathrm{Coker}}}
\def\Coim{\mathop{\mathrm{Coim}}}

% Boxtimes
%
\DeclareMathSymbol{\boxtimes}{\mathbin}{AMSa}{"02}

%
% Macros for moduli stacks/spaces
%
\def\QCohstack{\mathcal{QC}\!\mathit{oh}}
\def\Cohstack{\mathcal{C}\!\mathit{oh}}
\def\Spacesstack{\mathcal{S}\!\mathit{paces}}
\def\Quotfunctor{\mathrm{Quot}}
\def\Hilbfunctor{\mathrm{Hilb}}
\def\Curvesstack{\mathcal{C}\!\mathit{urves}}
\def\Polarizedstack{\mathcal{P}\!\mathit{olarized}}
\def\Complexesstack{\mathcal{C}\!\mathit{omplexes}}
% \Pic is the operator that assigns to X its picard group, usage \Pic(X)
% \Picardstack_{X/B} denotes the Picard stack of X over B
% \Picardfunctor_{X/B} denotes the Picard functor of X over B
\def\Pic{\mathop{\mathrm{Pic}}\nolimits}
\def\Picardstack{\mathcal{P}\!\mathit{ic}}
\def\Picardfunctor{\mathrm{Pic}}
\def\Deformationcategory{\mathcal{D}\!\mathit{ef}}
\usepackage{cleveref}
\crefname{originallemma}{original lemma}{original lemmas}
\Crefname{originallemma}{Original lemma}{Original lemmas}
\makeatletter
\newlabel{lemma-etale-stratified-finite}{{03S0}{source}{}{sourcepointer-1}{}}
\newlabel{lemma-etale-stratified-finite@cref}{{[originallemma][0][]03S0}{[1][0][]source}{}{}{}}
\makeatother
\crefname{originallemma}{original lemma}{original lemmas}
\Crefname{originallemma}{Original lemma}{Original lemmas}
\makeatletter
\newlabel{lemma-characterize-finite-locally-constant}{{03RV}{source}{}{sourcepointer-2}{}}
\newlabel{lemma-characterize-finite-locally-constant@cref}{{[originallemma][0][]03RV}{[1][0][]source}{}{}{}}
\makeatother
\crefname{originallemma}{original lemma}{original lemmas}
\Crefname{originallemma}{Original lemma}{Original lemmas}
\makeatletter
\newlabel{lemma-shriek-base-change}{{03S6}{source}{}{sourcepointer-3}{}}
\newlabel{lemma-shriek-base-change@cref}{{[originallemma][0][]03S6}{[1][0][]source}{}{}{}}
\makeatother
\crefname{originallemma}{original lemma}{original lemmas}
\Crefname{originallemma}{Original lemma}{Original lemmas}
\makeatletter
\newlabel{lemma-shriek-equals-star-finite-etale}{{03S7}{source}{}{sourcepointer-4}{}}
\newlabel{lemma-shriek-equals-star-finite-etale@cref}{{[originallemma][0][]03S7}{[1][0][]source}{}{}{}}
\makeatother
\crefname{originallemma}{original lemma}{original lemmas}
\Crefname{originallemma}{Original lemma}{Original lemmas}
\makeatletter
\newlabel{lemma-pushforward-locally-constant}{{095B}{source}{}{sourcepointer-5}{}}
\newlabel{lemma-pushforward-locally-constant@cref}{{[originallemma][0][]095B}{[1][0][]source}{}{}{}}
\makeatother
\title{Extension by zero of finite constant sheaves is constructible}
\author{The Stacks Project Authors}
\date{stacks-project:etale-cohomology.tex:local-snapshot-20261011; GNU FDL 1.2 or later}
\begin{document}
\maketitle
\noindent\textit{Editorial note.} All delivered mathematical text is the author's own, copied from the pinned Stacks Project chapter file etale-cohomology.tex; no statement or proof step was rewritten, reordered or removed.\par\medskip
\noindent\textit{Editorial note.} The proof is the author's complete argument for all three parts of the statement.\par\medskip
\noindent\textit{Editorial note.} Every labelled reference occurring inside the counted spans is carried in assembly.external\_references as a non-load-bearing navigation pointer (same-chapter ones at their definition site, cross-chapter ones at the invocation inside the counted span), so no load-bearing mathematics is replaced by a stub. No bibliography citation occurs inside the counted spans.\par\medskip
\noindent\textit{Editorial note.} The checked chapter list for same-name collisions is the set of stacks chapters staged in this repository: chow, cohomology, divisors, etale-cohomology, flat, more-algebra, more-morphisms, morphisms, spaces-more-morphisms, varieties together with this card's chapter etale-cohomology.\par\medskip
\noindent\textit{Editorial note.} The preamble is the chapter's own native preamble excerpt already used for this source: the Xy-pic setup (native preamble lines 14-18) and the native theorem declarations with the mathematical macros (native preamble lines 154-217).\par\medskip
\noindent\textit{Allowed standard prerequisites (editorial).}\begin{itemize}
\item Representable sheaf h\_U and extension by zero j\_! of constant sheaves. Same chapter, the sheaf represented by an etale scheme and the functor j\_! applied to constant sheaves. Supplies the sheaves whose constructibility is asserted.
\item Constructibility by a partition into constructible locally closed strata. Same chapter, the characterisation of constructibility through a partition and the result that an etale morphism of quasi-compact quasi-separated schemes is finite etale over the strata of a suitable partition. Supplies the partition with respect to which the sheaves are shown constructible.
\end{itemize}\medskip
\ifdefined\hypertarget\hypertarget{sourcepointer-1}{}\fi
\noindent\textit{Source reference.} Original lemma 03S0: sources/stacks-etale-cohomology/pinned-original-source.tex, source lines 12035--12035. Navigation pointer invoked at proof line 12199; the counted span contains the pointer only and the cited result is not reproduced inside it, so the pointer is registered as a non-load-bearing navigation stub in the established form for Stacks material. No bibliography citation is involved.\par\medskip
\ifdefined\hypertarget\hypertarget{sourcepointer-2}{}\fi
\noindent\textit{Source reference.} Original lemma 03RV: sources/stacks-etale-cohomology/pinned-original-source.tex, source lines 9994--9994. Navigation pointer invoked at proof line 12202; the counted span contains the pointer only and the cited result is not reproduced inside it, so the pointer is registered as a non-load-bearing navigation stub in the established form for Stacks material. No bibliography citation is involved.\par\medskip
\ifdefined\hypertarget\hypertarget{sourcepointer-3}{}\fi
\noindent\textit{Source reference.} Original lemma 03S6: sources/stacks-etale-cohomology/pinned-original-source.tex, source lines 11493--11493. Navigation pointer invoked at proof line 12209; the counted span contains the pointer only and the cited result is not reproduced inside it, so the pointer is registered as a non-load-bearing navigation stub in the established form for Stacks material. No bibliography citation is involved.\par\medskip
\ifdefined\hypertarget\hypertarget{sourcepointer-4}{}\fi
\noindent\textit{Source reference.} Original lemma 03S7: sources/stacks-etale-cohomology/pinned-original-source.tex, source lines 11571--11571. Navigation pointer invoked at proof line 12210; the counted span contains the pointer only and the cited result is not reproduced inside it, so the pointer is registered as a non-load-bearing navigation stub in the established form for Stacks material. No bibliography citation is involved.\par\medskip
\ifdefined\hypertarget\hypertarget{sourcepointer-5}{}\fi
\noindent\textit{Source reference.} Original lemma 095B: sources/stacks-etale-cohomology/pinned-original-source.tex, source lines 9973--9973. Navigation pointer invoked at proof line 12212; the counted span contains the pointer only and the cited result is not reproduced inside it, so the pointer is registered as a non-load-bearing navigation stub in the established form for Stacks material. No bibliography citation is involved.\par\medskip
\begin{lemma}
\label{lemma-jshriek-constructible}
Let $j : U \to X$ be an \'etale morphism of quasi-compact and
quasi-separated schemes.
\begin{enumerate}
\item The sheaf $h_U$ is a constructible sheaf of sets.
\item The sheaf $j_!\underline{M}$ is a constructible abelian sheaf
for a finite abelian group $M$.
\item If $\Lambda$ is a Noetherian ring and $M$ is a finite $\Lambda$-module,
then $j_!\underline{M}$ is a constructible sheaf of $\Lambda$-modules
on $X_\etale$.
\end{enumerate}
\end{lemma}

\begin{proof}
By Lemma \ref{lemma-etale-stratified-finite} there is a partition
$\coprod_i X_i$ such that $\pi_i : j^{-1}(X_i) \to X_i$ is finite \'etale.
The restriction of $h_U$ to $X_i$ is $h_{j^{-1}(X_i)}$ which is finite
locally constant by Lemma \ref{lemma-characterize-finite-locally-constant}.
For cases (2) and (3) we note that
$$
j_!(\underline{M})|_{X_i} =
\pi_{i!}(\underline{M}) =
\pi_{i*}(\underline{M})
$$
by Lemmas \ref{lemma-shriek-base-change} and
\ref{lemma-shriek-equals-star-finite-etale}.
Thus it suffices to show the lemma for $\pi : Y \to X$ finite \'etale.
This is Lemma \ref{lemma-pushforward-locally-constant}.
\end{proof}

\end{document}
